% ==================== 3.1 引言 ====================
\section{引言}
\label{sec:ch3_intro}

月面环境是影响巡视器自主行走的关键因素。准确、高效的月面环境数字孪生建模是实现巡视器自主行走的基础支撑。本章针对月面环境的特殊性和复杂性，系统研究月面地形、月壤环境、光照与温度环境的数字孪生建模方法。

月面环境具有以下显著特点：

\textbf{（1）地形复杂多变}

月球表面布满了撞击坑、山脉、峡谷、月海等地形特征。撞击坑的密度和大小分布广泛，从几厘米的小坑到数百公里的巨型盆地。地形坡度变化大，局部坡度可达30\degree 以上。这些地形特征对巡视器的行走路径选择、通过性评估、稳定性分析等具有重要影响。

\textbf{（2）月壤力学特性独特}

月壤是覆盖在月球基岩之上的松散物质层，平均厚度约5-10米。月壤颗粒棱角分明、粒径分布广，具有低内聚力、高内摩擦角的特点。月壤的力学参数随深度、位置变化显著，对巡视器车轮的沉陷和滑移行为有直接影响。

\textbf{（3）光照条件极端}

月球自转周期约为27.3天，一个昼夜长达约14个地球日。太阳高度角在一天内变化范围大，导致光照强度和阴影分布剧烈变化。极区存在永久阴影区和持续光照区，对巡视器的能源供给和热控设计至关重要。

\textbf{（4）温度环境恶劣}

月面昼夜温差可达300\degree C以上，白天最高温度约120\degree C，夜间最低温度约-170\degree C。温度的剧烈变化对巡视器的结构、电子设备和能源系统构成严峻挑战。

针对上述特点，本章分别研究月面地形数字孪生建模、月壤环境数字孪生建模、月面光照与温度环境建模，并通过实验验证模型的有效性。
